\documentclass[a4paper]{article}

\usepackage[utf8x]{inputenc}
\usepackage[T1]{fontenc}
\usepackage[francais,lang=FR,bloc,completemulti,ensemble,automarks]{automultiplechoice}
\usepackage{array}
\usepackage{tabularx}
\usepackage{multirow}
\usepackage{multicol}
\usepackage{enumitem}
\usepackage{courier}
\usepackage{listings}
\usepackage{color}
\usepackage{pdfpages}

\usepackage{tikz} %for oval
\usepackage[scaled]{helvet}

\definecolor{lightgray}{rgb}{.9,.9,.9}
\definecolor{darkgray}{rgb}{.2,.2,.2}
\definecolor{purple}{rgb}{0.65, 0.12, 0.82}

\renewcommand{\familydefault}{\sfdefault}

\lstdefinelanguage{JavaScript}{
  keywords={typeof, new, true, false, catch, function, return, null, catch, switch, var, if, in, while, for, do, else, case, break},
  keywordstyle=\color{blue}\bfseries,
  ndkeywords={class, export, boolean, throw, implements, import, this},
  ndkeywordstyle=\color{darkgray}\bfseries,
  identifierstyle=\color{black},
  sensitive=false,
  comment=[l]{//},
  morecomment=[s]{/*}{*/},
  commentstyle=\color{purple}\ttfamily,
  stringstyle=\color{darkgray}\ttfamily,
  morestring=[b]',
  morestring=[b]"
}

\lstset{
   language=JavaScript,
   %backgroundcolor=\color{lightgray},
   extendedchars=true,
   inputencoding=utf8x,
   literate={é}{{\'e}}1 {ê}{{\^{e}}}1 {è}{{\`{e}}}1 {à}{{\`{a}}}1,
   basicstyle=\footnotesize\ttfamily,
   showstringspaces=false,
   showspaces=false,
   numbers=left,
   numberstyle=\footnotesize,
   numbersep=9pt,
   tabsize=2,
   breaklines=true,
   showtabs=false,
   captionpos=b,
   escapeinside={(*@}{@*)}
}

\begin{document}

\AMCrandomseed{1237893}
\AMCboxColor{red}
\AMCboxDimensions{shape=oval}


%%% Scoring strategy:
\scoringDefaultM{e=0,v=0,b=0,m=0,p=0,d=0,mz=1}
%\scoringDefaultS{e=0,v=0,b=0,m=0,p=0,d=0,mz=1}

\def\AMCformQuestion#1{\vspace{\AMCformVSpace}\par {\sc Question #1 :} }

%%% liste des questions:
\input{questions_definition.tex}

\exemplaire{1}{

%%% debut de l'en-tête des copies :
\setlength{\parindent}{0pt}

\noindent{\textbf{Examen de fin de semestre} \hfill \includegraphics[]{src/graphics/logo_heg.png}}

\vspace{2ex}

\noindent{\textbf{\AMCUIclasses} \hfill \bf Session \AMCUIsession}

\vspace*{1cm}

{\renewcommand{\arraystretch}{4}
\begin{tabularx}{15cm}{p{6cm}p{0.5cm}p{7cm}}
\hline
\underline{Unité d'enseignement} & : & \textbf{\AMCUIcourse}    \\ \hline
\underline{Professeur-e/s} & : & \textbf{\AMCUIteacher}  \\ \hline
\underline{Moyens auxiliaires autorises} & : & \textbf{Toute documentation autorisée} \\ \hline
\underline{Date} & : & \textbf{\AMCUIdate} \\ \hline
\underline{Temps à disposition} & : & \textbf{\AMCUIduration} \\ \hline
\underline{Cotation utilisée pour la correction} & : & \textbf{\AMCUIpoints} \\ \hline
\underline{Consignes} & : & Ne répondez que sur la feuille de réponses! Aucune autre réponse ne sera prise en compte. Bon travail \\ \hline
\underline{Nom et prénom de l'étudiant-e} & : & \\ \hline
\end{tabularx}
}


\vspace*{2cm}Visa du professeur: \dotfill

\newpage

\vspace{3cm}

\begin{center}
	\textit{Cette page a été laissée intentionnellement blanche}
\end{center}


\newpage
%\multicolumn{1}{|c|}{\bf

{
\em\setlength{\parindent}{0pt}

Les questions faisant apparaître le symbole \multiSymbole{} peuvent
  présenter une ou plusieurs bonnes réponses. Les autres ont
  une unique bonne réponse.
\newline

A la fin de l'examen, vous rendrez cet énoncé et la feuille de réponses distribuée avec l'énoncé.
  Les deux documents doivent porter votre nom et prénom.
\newline


}\hspace*{\fill}
\vspace{1ex}

%%% fin de l'en-tête

\input{questions_layout.tex}

%\includepdf[pages=-]{src/Feuilles_de_brouillon_2015.pdf}

\AMCcleardoublepage

\AMCdebutFormulaire

%%% début de l'en-tête de la feuille de réponses
{\large\bf\noindent Feuille de réponses}

{\setlength{\parindent}{0pt}\hspace*{\fill}

%% \AMCcode{etu}{3}\hspace*{\fill}
\begin{minipage}[b]{6.5cm}\setlength{\parindent}{0pt}
%% $\longleftarrow{}$\hspace{0pt plus 1cm} veuillez ne rien coder ci-contre,
%% et inscrire votre nom et prénom \textit{lisiblement} ci-dessous.

\vspace{1ex}

\hfill\champnom{\fbox{
    \begin{minipage}{.9\linewidth}
      Nom:

      \vspace*{.5cm}\dotfill

      Prénom:

      \vspace*{.5cm}\dotfill

      \vspace*{1mm}
    \end{minipage}
  }}\hfill\vspace{1ex}

\end{minipage}\hspace*{\fill}

}

\vspace{1ex}

{\bf\em\setlength{\parindent}{0pt} Les réponses aux questions sont à donner \textit{exclusivement} sur cette feuille:
  les réponses données sur d'autres feuilles ne seront \textit{pas prises en compte}.

\vspace{1ex}
Coloriez complètement et proprement chaque case choisie, afin d'obtenir le meilleur contraste possible
par rapport à une case vide. Ne coloriez pas les cases au crayon à papier, n'utilisez pas de
couleur rouge (le programme correcteur ne voit pas les couleurs). Utilisez plutôt du noir ou du bleu foncé.
Pour corriger une éventuelle erreur de saisie, veuillez utiliser si possible un correcteur liquide (Tipp-Ex).
}

\vspace{1ex}

%%% fin de l'en-tête de la feuille de réponses

\begin{multicols}{2}
\setlength{\parindent}{0pt}
\formulaire
\end{multicols}


\vspace{20ex}

\fbox{%
\parbox{\textwidth}{
\textbf{Remarques à chaud}
\newline

\textit{Cet espace de liberté est mis à votre disposition pour toutes les remarques que vous désirez faire maintenant que le semestre et l'examen sont terminés... Le scanner n'arriveras pas à déchiffrer votre prose... il n'y a pas de risque que ça influence la correction!}

\vspace{9cm}
Bonnes vacances!
\newline
}}




\clearpage

}

\end{document}
